%==============================================================================
% FICHAMENTO CRÍTICO — Artificial Intelligence for Prediction of Dental
% Implant Success
% Conforme NBR 6023 (referências) e NBR 10520 (citações)
% Capítulo 19 — Predição de Sucesso de Implantes Dentários
%==============================================================================
\section{Fichamento: \citeonline{silva2024artificial}}

% --- 1. IDENTIFICAÇÃO ---
\subsection{Identificação}
\begin{tabular}{ll}
\textbf{Título:} & Artificial intelligence for prediction of dental implant
success \\
\textbf{Autor(es):} & Silva, Rafael; Oliveira, Paulo; Tanaka, Eiji;
Kim, Hyun-Jung; Schwartz, Marcos; Zembic, Anja \\
\textbf{Periódico:} & Oral Surgery, Oral Medicine, Oral Pathology and Oral Radiology \\
\textbf{Ano:} & 2024 \\
\textbf{Volume:} & 138 \\
\textbf{Número:} & 3 \\
\textbf{Páginas:} & 331--345 \\
\textbf{DOI:} & \url{https://doi.org/10.1016/j.oooo.2024.01.013} \\
\textbf{Qualis:} & A2 (Odontologia) \\
\textbf{Tipo:} & Artigo original com validação temporal prospectiva \\
\end{tabular}

% --- 2. OBJETIVO ---
\subsection{Objetivo}
Desenvolver e validar um modelo preditivo multimodal baseado em
aprendizado de máquina (ensemble stacking: Gradient Boosting, Random
Forest e Elastic Net) para predição individualizada de sucesso de
implantes dentários em 5 e 10 anos, integrando variáveis
demográficas (idade, sexo, tabagismo, diabetes), clínicas
(osseointegração primária, torque de inserção, ISQ, tipo de
protocolo de carga), radiográficas (densidade óssea alveolar medida
em CBCT, morfologia do rebordo) e do implante (marca, diâmetro,
comprimento, superfície, posição na arcada). O estudo busca superar
a abordagem dicotômica convencional (sucesso vs. falha) oferecendo
uma probabilidade contínua de sucesso com estimativa de tempo até
falha, permitindo planejamento terapêutico personalizado.

% --- 3. METODOLOGIA ---
\subsection{Metodologia}
\textbf{Desenho do estudo:} Estudo de coorte retrospectivo multicêntrico
com validação temporal prospectiva. \\
\textbf{Amostra:} 15.840 implantes instalados em 6.847 pacientes por
12 cirurgiões em 5 centros (Brasil, Suíça, Japão, Estados Unidos,
Alemanha), acompanhados por 5 a 12 anos. O desfecho primário foi
sucesso do implante (critérios de Albrektsson modificados: ausência
de mobilidade, perda óssea marginal < 1,5 mm no primeiro ano e < 0,2
mm/ano subsequente, ausência de sintomas persistentes). \\
\textbf{Métricas principais:} AUC-ROC, sensibilidade, especificidade,
valor preditivo positivo (VPP), valor preditivo negativo (VPN), F1-score,
calibração (curvas de calibração com intercepto e inclinação), Curva de
Decisão (DCA), Acurácia balanceada. \\
\textbf{Validação:} Validação cruzada 5-fold estratificada por centro;
validação temporal prospectiva com coorte independente de 1.240 implantes
instalados entre 2022--2023 com follow-up mínimo de 18 meses; análise de
subgrupos por densidade óssea (Cawood-Howell classes I--IV), regime de
carga (imediata, precoce, convencional) e setor da arcada (anterior,
pré-molar, molar).

% --- 4. RESULTADOS PRINCIPAIS ---
\subsection{Resultados Principais}
\begin{itemize}
  \item O ensemble stacking (Gradient Boosting + Random Forest + Elastic
        Net) alcançou AUC de 0,912 (IC 95\%: 0,891--0,933) na predição
        de sucesso em 5 anos na validação interna, superando modelos
        individuais (GB: 0,894, RF: 0,886, EN: 0,812) e regressão
        logística tradicional (AUC 0,743).
  \item A sensibilidade foi de 0,894 e a especificidade de 0,878 para
        o limiar otimizado por Youden, com VPP de 0,923 e VPN de 0,835,
        indicando que o modelo é particularmente confiável para
        confirmar o sucesso do implante (baixo falso-positivo).
  \item Na validação temporal prospectiva (n=1.240 implantes), a AUC foi
        de 0,883 (IC 95\%: 0,857--0,909), com degradação aceitável de
        2,9 pontos percentuais em relação à validação interna.
  \item As variáveis mais preditivas (SHAP) foram: densidade óssea à
        CBCT (gain médio 0,197), torque de inserção (0,154), ISQ na
        instalação (0,128), relação coroa-implante (0,112), tabagismo
        (0,098) e superfície do implante (0,074).
  \item O modelo identificou corretamente 72,4\% dos implantes que
        falharam nos primeiros 12 meses (falha precoce), com um ganho
        líquido de 0,31 na DCA para thresholds de probabilidade entre
        5\% e 30\%, região clinicamente relevante para decisão de
        protocolo de carga.
  \item A taxa de sucesso geral da amostra foi de 94,7\% em 5 anos e
        89,3\% em 10 anos, consistente com a literatura de
        implantodontia.
\end{itemize}

% --- 5. LIMITAÇÕES ---
\subsection{Limitações}
\begin{itemize}
  \item A definição de sucesso adotada (Albrektsson) é restritiva e não
        inclui desfechos relatados pelo paciente (PROMs) como
        satisfação estética, conforto mastigatório e impacto na
        qualidade de vida, que são relevantes para a tomada de decisão
        compartilhada.
  \item O follow-up da validação temporal (18 meses) é insuficiente para
        capturar falhas tardias (>5 anos) por peri-implantite, fratura
        do implante ou sobrecarga oclusal, limitando a validação do
        modelo para predição de longo prazo.
  \item A amostra contém predominantemente implantes de superfície
        tratada (87\%) com diâmetros de 3,5--4,5 mm, não
        representando adequadamente implantes de diâmetro estreito
        ($<$3,0 mm) ou superfície lisa usados em situações de rebordo
        atrófico, que apresentam taxas de sucesso distintas.
  \item O modelo não incorporou dados de carga oclusal quantitativa
        (forças mastigatórias, parafuncionamento) nem de qualidade
        óssea histomorfométrica, que são preditores conhecidos de
        sucesso a longo prazo.
  \item A classificação da densidade óssea por CBCT (escala
        Cawood-Howell adaptada) é semiquantitativa e apresenta
        discordância interexaminador moderada (kappa 0,62), que pode
        introduzir ruído na variável preditora mais importante do modelo.
\end{itemize}

% --- 6. CONTRIBUIÇÃO PARA O LIVRO ---
\subsection{Contribuição para o Livro}
Este artigo é a referência central do Capítulo 19 na seção de predição
de sucesso de implantes, fornecendo um estudo de caso completo de modelo
preditivo multimodal em implantodontia. O ensemble stacking é detalhado
no livro como exemplo de técnica avançada de combinação de modelos para
maximizar desempenho preditivo. A análise SHAP com variáveis clínicas
e radiográficas é utilizada para demonstrar como modelos de IA podem
quantificar o peso de fatores de risco conhecidos (tabagismo, densidade
óssea) e revelar interações não óbvias (torque $\times$ superfície). A
integração de critérios de sucesso baseados em evidência (Albrektsson)
com métricas de machine learning estabelece a ponte entre a prática
clínica e a modelagem computacional. Os thresholds de probabilidade da
DCA (5--30\%) são adotados no livro para definir zonas de decisão
clínica.

% --- 7. ANÁLISE CRÍTICA ---
\subsection{Análise Crítica}
\begin{itemize}
  \item \textbf{Validade interna:} Alta -- Ensemble com validação cruzada
        estratificada por centro e repetição 5 vezes; comparação com
        regressão logística como baseline mínimo; otimização de
        hiperparâmetros com busca em grade.
  \item \textbf{Validade externa:} Média-Alta -- Validação temporal
        prospectiva com coorte independente, embora com follow-up
        limitado (18 meses) que não cobre falhas tardias; centros
        multicêntricos em 5 países com diferentes protocolos
        cirúrgicos e restaurativos.
  \item \textbf{Reprodutibilidade:} Média -- A metodologia ensemble é
        descrita em detalhes e os hiperparâmetros ótimos são
        reportados, mas o dataset não é público (dados de pacientes)
        e o código não foi disponibilizado; os autores comprometeram-se
        a compartilhar o pipeline mediante solicitação acadêmica.
  \item \textbf{Relevância clínica:} Alta -- Predizer sucesso de implante
        com AUC 0,912 pode auxiliar na seleção de casos, planejamento
        de carga imediata e aconselhamento prognóstico, mas a
        aplicabilidade clínica depende da disponibilidade dos dados
        de entrada (especialmente ISQ e torque de inserção) na rotina
        do cirurgião-dentista.
  \item \textbf{Conflito de interesses:} Sim -- O estudo recebeu
        financiamento parcial de fabricante de implantes
        (Straumann/Institut Straumann AG) e um dos autores (Zembic, A.)
        é funcionária da empresa. O contrato de financiamento
        explicitamente garantiu independência editorial dos autores
        acadêmicos.
  \item \textbf{Viés identificado:} Viés de financiamento -- fabricante
        de implantes pode ter interesse na validação de modelos que
        favoreçam seus produtos (implantes de superfície tratada da
        marca), embora a análise SHAP não tenha identificado a marca
        do implante como preditora relevante; viés de verificação --
        implantes considerados sucesso podem apresentar peri-implantite
        subclínica não detectada pelos critérios de Albrektsson,
        superestimando a taxa de sucesso.
\end{itemize}

% --- 8. CITAÇÕES RELEVANTES ---
\subsection{Citações Relevantes}
\begin{citacao}
``The ensemble stacking model achieved an AUC of 0.912 (95\% CI
0.891--0.933) for 5-year implant success prediction. SHAP analysis
identified bone density on CBCT (mean gain 0.197), insertion torque
(0.154), and ISQ at placement (0.128) as the top predictors. The model
correctly identified 72.4\% of early implant failures (within 12
months), with a net benefit of 0.31 in decision curve analysis for
probability thresholds between 5\% and 30\%.'' (p. 339).
\end{citacao}

% --- 9. MAPEAMENTO SDD ---
\subsection{Mapeamento SDD}
\textbf{Problema:} Predição individualizada de sucesso de implantes
dentários em 5 e 10 anos com base em variáveis demográficas, clínicas,
radiográficas e do implante, superando a baixa acurácia (<0,75) da
regressão logística tradicional que assume linearidade e independência
entre preditores. \\
\textbf{Entrada:} 28 variáveis (demográficas: idade, sexo, tabagismo,
diabetes; clínicas: torque, ISQ, protocolo de carga; radiográficas:
densidade óssea CBCT, morfologia do rebordo; do implante: diâmetro,
comprimento, superfície, posição) de 15.840 implantes em 6.847
pacientes. \\
\textbf{Saída:} Probabilidade de sucesso em 5 e 10 anos com IC 95\%;
tempo estimado até falha (para implantes classificados como risco);
rank SHAP de fatores de risco individuais. \\
\textbf{Critérios:} AUC $\geq$ 0,90 na validação interna; degradação
$\leq$ 5\% na validação temporal; calibração com intercepto $\leq$ 0,05
e inclinação 0,90--1,10; VPP $\geq$ 0,90; ganho líquido positivo na DCA
para thresholds 5--30\%. \\
\textbf{Confiança:} 0,80 -- Estudo com amostra robusta e validação
temporal prospectiva, mas limitado por conflito de interesses com
fabricante de implantes, follow-up insuficiente para validação de
longo prazo (\geq5 anos) e ausência de PROMs como desfecho secundário.

% --- 10. REFERÊNCIA ABNT ---
\subsection{Referência ABNT}
\begin{verbatim}
SILVA, Rafael et al. Artificial intelligence for prediction of dental
implant success. Oral Surgery, Oral Medicine, Oral Pathology and Oral
Radiology, New York, v. 138, n. 3, p. 331-345, 2024. DOI:
10.1016/j.oooo.2024.01.013.
\end{verbatim}
\endinput
